% Created 2025-08-21 Thu 13:07
% Intended LaTeX compiler: pdflatex
\documentclass[11pt]{article}
\usepackage[utf8]{inputenc}
\usepackage[T1]{fontenc}
\usepackage{graphicx}
\usepackage{longtable}
\usepackage{wrapfig}
\usepackage{rotating}
\usepackage[normalem]{ulem}
\usepackage{amsmath}
\usepackage{amssymb}
\usepackage{capt-of}
\usepackage{hyperref}
\date{\today}
\title{}
\hypersetup{
 pdfauthor={},
 pdftitle={},
 pdfkeywords={},
 pdfsubject={},
 pdfcreator={Emacs 30.2 (Org mode 9.7.11)}, 
 pdflang={English}}
\begin{document}

\tableofcontents

\section{Single cell pre-analysis}
\label{sec:org15330a9}
\subsection{Outline}
\label{sec:org66c1e9f}
\begin{itemize}
\item filtering cell by cell status, to remove bad quality cluster.
\item RNA, ATAC, Joint embedding.
\item pick a resolution to select enough number of cluster
\end{itemize}
\subsection{Draft}
\label{sec:org0eccdee}
In the single-cell pre-analysis, we first filtered cells based on their status to remove low-quality clusters.
Next, we performed RNA and ATAC joint embedding to integrate the two modalities. calculate the PCA then use harmony to remove the batch effect
for each modality and then concatenate the embedding results to create a joint embedding.
We then selected a resolution that allowed us to identify a sufficient number of clusters for further analysis.
\section{Peak calling}
\label{sec:org08ba434}
\subsection{Outline}
\label{sec:orga4354f3}
\begin{itemize}
\item Separate fragment to each cluster
\item macs2 recalling
\item generate cluster peaks count and union peaks count
\end{itemize}
\subsection{Draft}
\label{sec:org3bc4059}
In the peak recalling step, we separated the ATAC-seq fragments by cluster and applied MACS2 to call peaks for each cluster.
\section{Network analysis}
\label{sec:orga3221d6}
\subsection{Outline}
\label{sec:orgfa4473e}
\begin{itemize}
\item In each line calculate logFC for each gene between Coc and Suc group
\end{itemize}
\subsubsection{LogFC}
\label{sec:orgf1b6c2c}
\begin{itemize}
\item Calculate correlation of logFC between each line for check the similarity in each line.
\item Separate similarity to two group Within group (correlation between same preference) and 
Between group (correlation between different preference)
\item Do test to check the significance of the difference between the two groups
\end{itemize}
\subsubsection{Differential TF activity}
\label{sec:orgfb56e89}
\begin{itemize}
\item Calculate DTFA score for each TF in each line, where DTFA score is weighted sum of logFC of each gene regulated by the TF.
\item Calculate correlation of DTFA score between each line for check the similarity in each line.
\item Separate similarity to two group Within group (correlation between same preference) and 
Between group (correlation between different preference)
\item Do test to check the significance of the difference between the two groups
\end{itemize}
\subsubsection{Significant TF}
\label{sec:org5fcac7b}
\begin{itemize}
\item Calculate DTFA and Normalize it
\item For each TF, compare DTFA between mid-preference and high-preference in male and female separately.
\end{itemize}
\subsubsection{fisher z-based test}
\label{sec:orgca29950}
\subsubsection{Correlation analysis using Fisher z test}
\label{sec:orgeaf01bb}
To assess whether correlations were systematically stronger within groups
than between groups, we applied a Fisher z–based test. For each pair of
features, Pearson correlations were computed and annotated according to
whether both belonged to the same preference group (within-group) or to
different groups (between-group). Correlation coefficients were transformed
using the Fisher z-transformation,
\(z = \tfrac{1}{2}\ln\big(\tfrac{1+r}{1-r}\big)\), which stabilizes
variance and approximates normality.

We then calculated the mean Fisher z for within-group and between-group
edges and obtained their difference (\(\mu\)). The standard error of this
difference was estimated as

\(sd = \sqrt{\frac{1}{N-3}\left(\frac{1}{n_{\text{within}}} +
\frac{1}{n_{\text{between}}}\right)}\),

where \(N\) is the number of samples used to compute each correlation, and
\(n_{\text{within}}\), \(n_{\text{between}}\) are the counts of within- and
between-group edges, respectively.

The resulting test statistic was \(Z = \mu / sd\). A one-tailed p-value
was derived as \(p = 1 - \Phi(Z)\), where \(\Phi\) is the cumulative
distribution function of the standard normal distribution. This tests the
hypothesis that within-group correlations are greater than between-group
correlations.
\subsection{Draft}
\label{sec:orgafa856e}
For each genetic line, we calculated gene-level log fold change (logFC) between the cocaine (Coc) and sucrose (Suc) groups.
To assess similarity of treatment effects across lines, we computed correlations of logFC values for each pair of lines.
These correlations were grouped into within-group (same preference) and between-group (different preference) categories.
Statistical testing was performed to evaluate whether within-group correlations were significantly higher than between-group correlations.
In this analysis we found that there are no significant differences between the two groups.
This suggests that the treatment effects on gene expression are similar across lines, regardless of their preference for cocaine or sucrose.

For Gene regulatory network(GRN) analysis, we computed a Differential TF Activity (DTFA) score for each TF in each line.
The DTFA score is a weighted sum of logFC values of its target genes.
Similar to gene-level analysis, we assessed the pairwise correlation of DTFA scores across lines,
separated them into within- and between-group categories, and tested for significant differences.
In the analysis of DTFA score, we found that there are significant differences between the two groups,
indicating that TF activity patterns differ more between preferences than within them.

For significant TFs, we compared DTFA scores between mid-preference and high-preference, for the cell types that we found significant differences in the previous analysis.
Each cell types have their own set of significant TFs.
\end{document}
